\scanpage{291}
\uline{Structure de certains groupes divisibles}.

\deleted{\unsure{Le|On}} Les \deleted{$p$}groupes de torsion divisibles $\deleted{\unsure{G|A}}M$ sont ceux qui sont isomorphes à des sommes directes de groupes $\mathbb{Q}_p/\mathbb{Z}_p$ ($p$ premier); Ceux qui sont $p$-primaires sont ceux isom. à des \deleted{sommes directes de} groupes $(\mathbb{Q}_p/\mathbb{Z}_p)^{(I)}$, où le cardinal de $I$ est bien déterminé; c'est aussi la dimension de l'espace vectoriel sur $p$ $\,{}_p\deleted{G}M=\operatorname{Hom}(\mathbb{Z}/p\mathbb{Z},\deleted{G}M)$.

\deleted{Les groupes de torsion $p$-primaires sont d'ailleurs isomorphes}\inserted{Soit $\deleted{G}M$ un groupe.}

Conditions équivalentes (i) $\deleted{G}M$ groupe de torsion $p$-primaire, et $\,{}_p\deleted{G}M$ est fini. (ii) $\deleted{G}M$ isomorphe à un sous-groupe d'un $(\mathbb{Q}_p/\mathbb{Z}_p)^{\unsure{N|I|n}}$, \deleted{\unsure{$I$ ens. de c|avec $n$}} un entier $\geq0$. (iii) $\deleted{G}M$ isomorphe à un groupe $F+(Q_p/\mathbb{Z}_p)^\rho$, où $F$ est un $p$-groupe de torsion fini, et $\rho$ un entier fini.

% slanted
\margin{modules sur un anneau de \unsure{Dedekind|Dieudonné}.}

Sous ces conditions, $\rho$ est défini de façon unique [c'est aussi $\operatorname{rang}_{\mathbb{Z}/p\mathbb{Z}}\,{}_p\deleted{G}M^0$, où $M^0$ est le plus grand sous-groupe divisible de $\deleted{G}M$] ainsi que la structure de $F$ [c'est $M/\deleted{G}M^0$].

On peut de deux façons différentes faire correspondre à $M$ un module de type fini sur $\mathbb{Z}_p$. On peut poser

\[
T_p(\deleted{G}M)=\operatorname{Hom}(M,\mathbb{Q}_p/\mathbb{Z}_p)=\underset{n}{\varprojlim}\operatorname{Hom}({}_{p^n}\deleted{G}M,\mathbb{Q}_p/\mathbb{Z}_p)
\]

\margin{$T_p$ \uline{foncteur} \uline{contravariant} \uline{en} $M$}

[$M\text{ en }T_p(M)$ établit une \inserted{anti}équivalence de la catégorie des $M$ satisfaisant (i), et des modules de type fini sur $\mathbb{Z}_p$.

On peut prendre aussi

\[
T^p(M)=\varprojlim{}_{p^n}M
\]

% slanted
\margin{$T^p$ foncteur \uline{covariant} en $M$}

où pour $m>n$, on envoie $\,{}_{p^m}M$ dans $\,{}_{p^n}M$ par $p^{m-n}$.

% archive sheet number: 290

